\fig{em5:damping-time}{相同初电荷、零初电流时，三种阻尼的电荷时间曲线}{
\begin{tikzpicture}[x=.9cm,y=2cm]
\draw[axis](0,0)--(9.5,0)node[right]{$\omega_0t$};
\draw[axis](0,-.7)--(0,1.25)node[above]{$q/Q_0$};
\draw[world,blue,domain=0:9,samples=160]plot(\x,{exp(-.2*\x)*(cos(sqrt(.96)*\x r)+.2/sqrt(.96)*sin(sqrt(.96)*\x r))});
\draw[world,red,dashed,domain=0:9,samples=100]plot(\x,{(1+\x)*exp(-\x)});
\draw[world,black,dash dot,domain=0:9,samples=100]
 plot(\x,{((2+sqrt(3))*exp((-2+sqrt(3))*\x)-(2-sqrt(3))*exp((-2-sqrt(3))*\x))/(2*sqrt(3))});
\node[left]at(0,1){$1$};
\node[blue]at(1.3,-1){欠阻尼};
\node[red]at(4.4,-1){临界阻尼};
\node at(7.7,-1){过阻尼};
\end{tikzpicture}}

图\ref{em5:damping-time}统一使用无阻尼角频率$\omega_0=1/\sqrt{LC}$量度时间。三条曲线分别取阻尼比$R/(2\sqrt{L/C})$为$0.2,1,2$，起点高度和初始斜率相同。蓝线穿过横轴，表示电容换极；另外两线在本图初态下保持正电荷并渐近衰减。不同初电流可能改变过零情况，因而不能把这一特定初态的图推广为任意初态都不换极。临界与过阻尼的判定依据仍是特征根及轨道性质。
